\clearpage
\section{Actual adapted source
inventory}\label{d-actual-adapted-source-inventory}

Technical derivation for review. Equation and subsection numbers in this
appendix are local to the source derivation. Historical local labels are
retained, including numbering gaps or nonsequential labels. The
accompanying source notes retain the original expressions for
comparison.

\subsection{1. Notation and exact
origin}\label{d-notation-and-exact-origin}

Use the kinetic-one normalization and the intrinsic slice in Appendix A.
The only internal thinness hypothesis is

% DISPLAY d.1
% SOURCE alpha_a+alpha_b <= kappa r_ab,
\begin{gather*}
\begin{aligned}
& \alpha _{a}+\alpha _{b} \le  \kappa  r_{ab},
\end{aligned}\notag
\end{gather*}

with fixed sufficiently small \(\kappa\). Set \(r_{*}=\min  r_{ab}\).
All finite constants below depend on \(N\) but not on a ratio of
unrelated gaps.

To avoid overloading center \(C\), use these coefficient maps:

% DISPLAY d.2
% SOURCE F = E*B_(Z+A),             C = P_N B_A,
% SOURCE L = E*B_Y,                 N_Y = P_N B_Y,
% SOURCE I_Y = P_I B_Y,             C_Y = P_C B_Y,
% SOURCE K = R^-1 C_Y*,             J_2 = K C_Y.
\begin{gather*}
\begin{aligned}
& F = E^{*}B_{Z+A}, \quad  C = P_{N} B_{A}, \\
& L = E^{*}B_{Y}, \quad  N_{Y} = P_{N} B_{Y}, \\
& I_{Y} = P_{I} B_{Y}, \quad  C_{Y} = P_{C} B_{Y}, \\
& K = R^{-1} C_{Y}^{*}, \quad  J_{2} = K C_{Y}.
\end{aligned}\notag
\end{gather*}

Thus \(F\) is pair-block diagonal, \(C\) is pair-block diagonal as a map
from gauge parameters into normal fast coordinates, and

% DISPLAY d.3
% SOURCE M=F+L-J_2,   W=I+KK*.
\begin{gather*}
\begin{aligned}
& M=F+L-J_{2}, \quad  W=I+KK^{*}.
\end{aligned}\notag
\end{gather*}

The maps \(L\) and \(N_{Y}\) between distinct fast root pairs have
triangle support. They have no same-pair block. The maps \(I_{Y}\) and
\(C_{Y}\) use one root and its endpoints. The center metric is
\(g_{C}=I+K^{*}K\); the internal and fast tangent metrics are exactly
identity.

The definitive flattened kinetic form is

% DISPLAY d.4
% SOURCE ||(p_H+i rho)u||_(g_S^-1)^2 + ||D u||²,
% SOURCE D=W^(1/2)M^-*[S_f-O*Sg_S^-1(p_H+i rho)],
% SOURCE rho=d_H log d,       d=sqrt(det M).                     (1.1)
\begin{gather*}
\begin{aligned}
& \Vert (p_{H}+i \rho )u\Vert _{g_{S}^{-1}}^{2} + \Vert D u\Vert ^{2}, \\
& D=W^{1/2}M^{-*}[S_{f}-O^{*}Sg_{S}^{-1}(p_{H}+i \rho )], \\
& \rho =d_{H} \log  d, \quad  d=\sqrt{\det  M}.
\end{aligned}\tag{1.1}
\end{gather*}

All products below inherit this ordering. Fast adjoints are taken in the
flat fiber measure. The normal connection is retained in center
derivatives. Nothing is obtained by substituting a momentum constraint
into a squared operator.

\subsection{2. Exact coefficients needed through
H2}\label{d-exact-coefficients-needed-through-h2}

Use homogeneous weights

% DISPLAY d.5
% SOURCE (Z,A) -> L_scale (Z,A),  Y -> L_scale^-1/2 Y,
% SOURCE p_s -> L_scale^-1 p_s,  p_F -> L_scale^1/2 p_F.
\begin{gather*}
\begin{aligned}
& (Z,A) \to  L_{\mathrm{scale}} (Z,A), \quad  Y \to  L_{\mathrm{scale}}^{-1/2} Y, \\
& p_{s} \to  L_{\mathrm{scale}}^{-1} p_{s}, \quad  p_{F} \to  L_{\mathrm{scale}}^{1/2} p_{F}.
\end{aligned}\notag
\end{gather*}

Here \(s\) comprises centers and all internal coordinates. Keep the
internal potential and internal Yukawa terms in their complete internal
Hamiltonians throughout.

Write

% DISPLAY d.6
% SOURCE log d = log d_0 + phi_2 + terms of Y-degree >=3,
% SOURCE d_0=sqrt(det F),
% SOURCE phi_2=-(1/2)tr(F^-1 J_2)
% SOURCE         -(1/4)tr(F^-1 L F^-1 L),                       (2.1)
% SOURCE rho_0=d_s log d_0,    rho_F2=d_F phi_2,
% SOURCE u_s=p_s+i rho_0.
\begin{gather*}
\begin{aligned}
& \log  d = \log  d_{0} + \phi _{2} + \text{ terms of Y-degree } \ge 3, \\
& d_{0}=\sqrt{\det  F}, \\
& \phi _{2}=-(1/2)\operatorname{tr} (F^{-1} J_{2}) \\
& -(1/4)\operatorname{tr} (F^{-1} L F^{-1} L),
\end{aligned}\tag{2.1} \\
\begin{aligned}
& \rho _{0}=d_{s} \log  d_{0}, \quad  \rho _{F2}=d_{F} \phi _{2}, \\
& u_{s}=p_{s}+i \rho _{0}.
\end{aligned}\notag
\end{gather*}

There is no degree-one term because \(\operatorname{tr} (F^{-1} L)=0\)
for every A, not only at \(A=0\). The triangle term in (2.1) must be
retained.

The Gauss derivative stack has coefficients

% DISPLAY d.7
% SOURCE D = D_0 + D_1 + D_2 + higher homogeneous orders,
\begin{gather*}
\begin{aligned}
& D = D_{0} + D_{1} + D_{2} + \text{ higher homogeneous orders },
\end{aligned}\notag
\end{gather*}

of scaling degrees \(1/2\), \(-1\), and \(-5/2\), respectively. They are

% DISPLAY d.8
% SOURCE D_0 = -F^-1 C* p_F,                                     (2.2)
% SOURCE 
% SOURCE D_1 = F^-1(S_f-N_Y* p_F)
% SOURCE          +F^-1 L* F^-1 C* p_F,                         (2.3)
\begin{gather*}
\begin{aligned}
& D_{0} = -F^{-1} C^{*} p_{F},
\end{aligned}\tag{2.2} \\
\begin{aligned}
& D_{1} = F^{-1}(S_{f}-N_{Y}^{*} p_{F}) \\
& +F^{-1} L^{*} F^{-1} C^{*} p_{F},
\end{aligned}\tag{2.3}
\end{gather*}

and

% DISPLAY d.9
% SOURCE D_2 = -F^-1[I_Y* u_I+(C_Y*+F K)u_C+i C* rho_F2]
% SOURCE       -F^-1 L* F^-1(S_f-N_Y* p_F)
% SOURCE       -F^-1 L* F^-1 L* F^-1 C* p_F
% SOURCE       -F^-1 J_2* F^-1 C* p_F
% SOURCE       -(1/2)KK* F^-1 C* p_F.                           (2.4)
\begin{gather*}
\begin{aligned}
& D_{2} = -F^{-1}[I_{Y}^{*} u_{I}+(C_{Y}^{*}+F K)u_{C}+i C^{*} \rho _{F2}] \\
& -F^{-1} L^{*} F^{-1}(S_{f}-N_{Y}^{*} p_{F}) \\
& -F^{-1} L^{*} F^{-1} L^{*} F^{-1} C^{*} p_{F} \\
& -F^{-1} J_{2}^{*} F^{-1} C^{*} p_{F} \\
& -(1/2)KK^{*} F^{-1} C^{*} p_{F}.
\end{aligned}\tag{2.4}
\end{gather*}

These follow by expanding

% DISPLAY d.10
% SOURCE M^-*=F^-1-F^-1L*F^-1
% SOURCE          +F^-1L*F^-1L*F^-1+F^-1J_2*F^-1+...,
% SOURCE W^(1/2)=I+(1/2)KK*+...,
% SOURCE O*Sg_S^-1p_H
% SOURCE    =C*p_F+N_Y*p_F+I_Y*p_I+(C_Y*+F K)p_C+... .          (2.5)
\begin{gather*}
\begin{aligned}
& M^{-*}=F^{-1}-F^{-1}L^{*}F^{-1} \\
& +F^{-1}L^{*}F^{-1}L^{*}F^{-1}+F^{-1}J_{2}^{*}F^{-1}+\ldots , \\
& W^{1/2}=I+(1/2)KK^{*}+\ldots , \\
& O^{*}Sg_{S}^{-1}p_{H} \\
& =C^{*}p_{F}+N_{Y}^{*}p_{F}+I_{Y}^{*}p_{I}+(C_{Y}^{*}+F K)p_{C}+\ldots  .
\end{aligned}\tag{2.5}
\end{gather*}

For example \(L^{*}K p_{C}\) first enters the next Gauss order, and
\(g_{C}^{-1}-I\) first contributes to direct center kinetic energy at
scaling degree -5; neither belongs to H2. The \(F K\) in (2.4) is
essential and becomes \(C_{Y}^{*}\) at \(A=0\), so the centered
center-orbital coefficient is twice \(C_{Y}^{*}\).

The exact leading fast Hamiltonian comprises the direct \(p_{F}^{2}\),
\(D_{0}^{*}D_{0}\), the full adapted quadratic potential, and the fast
fermion mass. Its product vacuum is \(U(A)\), with energy
\(E_{0,\mathrm{total}}\). Define

% DISPLAY d.11
% SOURCE H_exc = H_fast-E_0,total.
\begin{gather*}
\begin{aligned}
& H_{\mathrm{exc}} = H_{\mathrm{fast}}-E_{0,\mathrm{total}}.
\end{aligned}\notag
\end{gather*}

The vacuum-energy multiplier is handled by Appendix \(B\). It is not
inserted into a virtual excitation denominator.

The next two homogeneous form coefficients are exactly

% DISPLAY d.12
% SOURCE H_1 = V_3+Y_Y +D_0*D_1+D_1*D_0,                       (2.6)
% SOURCE 
% SOURCE H_2 = ||u_s (.)||² +V_4
% SOURCE         +2 Re <p_F(.),i rho_F2(.)>
% SOURCE         +||D_1(.)||²+2 Re <D_0(.),D_2(.)>.             (2.7)
\begin{gather*}
\begin{aligned}
& H_{1} = V_{3}+Y_{Y} +D_{0}^{*}D_{1}+D_{1}^{*}D_{0},
\end{aligned}\tag{2.6} \\
\begin{aligned}
& H_{2} = \Vert u_{s} (.)\Vert ^{2} +V_{4} \\
& +2 \operatorname{Re}  \langle p_{F}(.),i \rho _{F2}(.)\rangle \\
& +\Vert D_{1}(.)\Vert ^{2}+2 \operatorname{Re}  \langle D_{0}(.),D_{2}(.)\rangle .
\end{aligned}\tag{2.7}
\end{gather*}

Here \(Y_{Y}\) is the Yukawa term linear in \(Y\), including both
slow-fast and fast-fast fermion terms. \(V_{3}\) is the cross term
between \([Z+A,Y]\) and \([Y,Y]\); \(V_{4}\) is the pure quartic fast
potential. The internal potential term paired with \([Y,Y]\) is already
in the adapted quadratic potential. Equations (2.6)-(2.7) therefore
include the terms which appear as anticommutators of the internal
potential charge with higher inverse-Gauss charge coefficients; there is
no permission to omit them in a charge derivation. This Hamiltonian
derivation fixes them directly from the original scalar potential and
the exact kinetic form.

\subsection{3. Actual H1 source and its graph
weights}\label{d-actual-h1-source-and-its-graph-weights}

Define the actual leading source

% DISPLAY d.13
% SOURCE S(A) = H_1 U(A).                                        (3.1)
\begin{gather*}
\begin{aligned}
& S(A) = H_{1} U(A).
\end{aligned}\tag{3.1}
\end{gather*}

It is complementary by boson parity. Each summand is a finite
Gaussian-polynomial/Fock map on the full slow Clifford module. It has
the graph decomposition

% DISPLAY d.14
% SOURCE S=sum_e S_e+sum_t S_t,
\begin{gather*}
\begin{aligned}
& S=\sum _{e} S_{e}+\sum _{t} S_{t},
\end{aligned}\notag
\end{gather*}

where an edge source has exactly two excitations on that edge
(\(N_{e}=2\)) and a triangle source has exactly one excitation on every
triangle edge (\(N_{e}=N_{f}=N_{g}=1\)). In the pair case these are one
bosonic and one fast-fermionic excitation. Graph subspaces remain
orthogonal and invariant under the adapted quadratic oscillator. Parity
alone would not be the complete finite-sector statement.

The inventory is explicit:

\begin{itemize}
\tightlist
\item
  \(Y_{Y} U\): a pair has \(Y_{e} \theta _{e} \theta _{s}\); a triangle
  has \(Y_{e} \theta _{f} \theta _{g}\).
\item
  \(V_{3} U\): only three distinct triangle edges, one \(Y\) on each.
\item
  \(D_{0}^{*}F^{-1} S_{f} U\) and its adjoint: pair terms have one
  \(p_{e}\) and one fast fermion; triangle terms have
  \(p_{e} \theta _{f} \theta _{g}\).
\item
  \(D_{0}^{*}F^{-1} N_{Y}^{*}p_{F} U\) and its adjoint: a triangle has
  \(p_{e} Y_{f} p_{g}\).
\item
  The last term of \(D_{1}\) yields the same triangle support with one
  additional factor \(F_{f}^{-1} C_{f}^{*}\).
\end{itemize}

A slow internal derivative does not occur in H1. Only the pair source
contains internal/center Clifford generators; the triangle source acts
trivially on the internal Clifford factors.

For a triangle let \(s_{t}=\sum _{e \in  t}r_{e}\),
\(w_{t}=\sum _{e \in  t}r_{e}^{-2}\). Pairwise thinness gives

% DISPLAY d.15
% SOURCE ||F_e^-1|| <= C/r_e,    ||F_e^-1 C_e*|| <= C kappa.
\begin{gather*}
\begin{aligned}
& \Vert F_{e}^{-1}\Vert  \le  C/r_{e}, \quad  \Vert F_{e}^{-1} C_{e}^{*}\Vert  \le  C \kappa .
\end{aligned}\notag
\end{gather*}

On a fixed finite Gaussian degree, \(Y_{e}\) costs \(C r_{e}^{-1/2}\),
and \(p_{e}\) costs \(C r_{e}^{1/2}\). Thus

% DISPLAY d.16
% SOURCE ||S_e||² <= C_N/r_e,
% SOURCE ||S_t||² <= C_N s_t w_t.                               (3.2)
\begin{gather*}
\begin{aligned}
& \Vert S_{e}\Vert ^{2} \le  C_{N}/r_{e}, \\
& \Vert S_{t}\Vert ^{2} \le  C_{N} s_{t} w_{t}.
\end{aligned}\tag{3.2}
\end{gather*}

For the cubic potential, use

% DISPLAY d.17
% SOURCE s_t/(r_e r_f r_g)
% SOURCE    =1/(r_e r_f)+1/(r_f r_g)+1/(r_g r_e) <= w_t.
\begin{gather*}
\begin{aligned}
& s_{t}/(r_{e} r_{f} r_{g}) \\
& =1/(r_{e} r_{f})+1/(r_{f} r_{g})+1/(r_{g} r_{e}) \le  w_{t}.
\end{aligned}\notag
\end{gather*}

For the Gauss triangle \(p_{e} Y_{f} p_{g}/r_{e}\), its norm is bounded
by

% DISPLAY d.18
% SOURCE C kappa sqrt(r_g/(r_e r_f))
% SOURCE    <= C kappa sqrt(s_t) (r_e r_f)^-1/2.
\begin{gather*}
\begin{aligned}
& C \kappa  \sqrt{r_{g}/(r_{e} r_{f})} \\
& \le  C \kappa  \sqrt{s_{t}} (r_{e} r_{f})^{-1/2}.
\end{aligned}\notag
\end{gather*}

The inverse-L correction has one additional bounded factor \(\kappa\).
These are shape-independent estimates, including when one edge is much
shorter than the other two.

Transporting only the finite oscillator coefficient spaces to their
centered spaces by the actual pair functional calculus gives the
coefficientwise perturbation estimates

% DISPLAY d.19
% SOURCE ||S_e(A)-S_e(0)|| <= C_N kappa r_e^-1/2,
% SOURCE ||S_t(A)-S_t(0)|| <= C_N kappa sqrt(s_t w_t).             (3.3)
\begin{gather*}
\begin{aligned}
& \Vert S_{e}(A)-S_{e}(0)\Vert  \le  C_{N} \kappa  r_{e}^{-1/2}, \\
& \Vert S_{t}(A)-S_{t}(0)\Vert  \le  C_{N} \kappa  \sqrt{s_{t} w_{t}}.
\end{aligned}\tag{3.3}
\end{gather*}

This is an estimate for the actual source (3.1), not a definition by
transport. Every new Gauss H1 term has an explicit factor
\(F^{-1} C^{*}\), and the remaining actual coefficients have relative
\(O(\kappa )\) variation in the pair matrices and vacuum factors.

\subsection{4. Compressed H2 and its actual scalar/Clifford
multiplier}\label{d-compressed-h2-and-its-actual-scalarclifford-multiplier}

Separate the complete slow Berry-covariant form

% DISPLAY d.20
% SOURCE H_s^B=T_cent^B+sum_a(T_a^B+V_a+Y_a)
\begin{gather*}
\begin{aligned}
& H_{s}^{B}=T_{\mathrm{cent}}^{B}+\sum _{a}(T_{a}^{B}+V_{a}+Y_{a})
\end{aligned}\notag
\end{gather*}

from the vacuum compression of (2.7). The remainder is a multiplication
operator \(W(A)\) on the entire slow Clifford fiber. The following
procedure is an exact specification of \(W\):

\begin{enumerate}
\def\labelenumi{\arabic{enumi}.}
\item
  In \(\Vert u_{s}(Uf)\Vert ^{2}\), retain
  \(\Vert \nabla _{s}^{B} f\Vert ^{2}\); put the Born-Huang multiplier
  and \(\vert \rho _{0}\vert ^{2}+\operatorname{div}  \rho _{0}\) into
  \(W\).
\item
  Put the vacuum expectations of \(V_{4}\), the displayed fast-density
  cross term, and \(D_{1}^{*}D_{1}\) into \(W\).
\item
  In \(2 \operatorname{Re} \langle D_{0}(Uf),D_{2}(Uf)\rangle\), keep
  the written ordering. For the slow-orbital part put

  % DISPLAY d.21
  % SOURCE L_mu = component_mu{F^-1[I_Y*,C_Y*+F K]},
  % SOURCE <D_0 U,-L_mu U> = i j_mu,       j_mu real.            (4.1)
  \begin{gather*}
  \begin{aligned}
  & L_{\mu} = \operatorname{component} _{\mu}\{F^{-1}[I_{Y}^{*},C_{Y}^{*}+F K]\}, \\
  & \langle D_{0} U,-L_{\mu} U\rangle  = i j_{\mu}, \quad  j_{\mu} \text{ real }.
  \end{aligned}\tag{4.1}
  \end{gather*}

  This coefficient is purely imaginary: the bosonic Gaussian and the
  orbital coefficient maps are real, while \(D_{0}\) contains one fast
  momentum. The fermionic factor is its norm. Its first-order
  coefficient on \(f\) therefore integrates to the scalar multiplier
  \(-\operatorname{div}  j\). All remaining terms, including derivatives
  on \(U\) and \(i \rho _{0}\), are placed in \(W\). A scalar Berry term
  in \(p_{\mu} U\) has zero real pairing with \(i j_{\mu}\), so this
  step is gauge covariant.
\end{enumerate}

An equivalent exact scalar formula removes even the implicit derivative
bookkeeping. Write \(U=g(Y;s)v(s)\), with \(g\) real, put

% DISPLAY d.22
% SOURCE A_op=F^-1 C*,    a=A_op d_F log g,
% SOURCE B_mu=component_mu{F^-1[I_Y*,C_Y*+F K]}.
\begin{gather*}
\begin{aligned}
& A_{\mathrm{op}}=F^{-1} C^{*}, \quad  a=A_{\mathrm{op}} d_{F} \log  g, \\
& B_{\mu}=\operatorname{component} _{\mu}\{F^{-1}[I_{Y}^{*},C_{Y}^{*}+F K]\}.
\end{aligned}\notag
\end{gather*}

Then \(D_{0} U=i a U\), and the entire slow-orbital compression is

% DISPLAY d.23
% SOURCE V_orb=-sum_mu <partial_mu(a dot B_mu)>_U
% SOURCE          -2 sum_mu rho_(0,mu)<a dot B_mu>_U.            (4.3)
\begin{gather*}
\begin{aligned}
& V_{\mathrm{orb}}=-\sum _{\mu} \langle \partial _{\mu}(a \cdot  B_{\mu})\rangle _{U} \\
& -2 \sum _{\mu} \rho _{0,\mu}\langle a \cdot  B_{\mu}\rangle _{U}.
\end{aligned}\tag{4.3}
\end{gather*}

The derivative acts on the displayed fast coefficient; Gaussian-density
differentiation has already canceled the derivative of the fiber
expectation. Fock/Berry derivative terms have zero real part. The direct
fast-density term and the \(\rho _{F2}\) term in \(D_{2}\) combine
exactly to

% DISPLAY d.24
% SOURCE V_fast-density=tr_F[(I+C F^-2 C*) Hess_Y phi_2].          (4.4)
\begin{gather*}
\begin{aligned}
& V_{\mathrm{fast\text{-}density}}=\operatorname{tr} _{F}[(I+C F^{-2} C^{*}) \operatorname{Hess} _{Y} \phi _{2}].
\end{aligned}\tag{4.4}
\end{gather*}

Since \(\phi _{2}\) is quadratic in \(Y\), its Hessian is a coefficient
matrix independent of \(Y\). These formulas follow by one integration by
parts and retain all ordering terms.

No slow derivative has been discarded or bounded by a momentum cutoff.
In particular differentiating \(j\), or the coefficient in (4.3),
includes internal derivatives at \(A=0\). Those terms participate in the
centered cancellation.

\subsubsection{4.1 Complete root support and coefficient
bounds}\label{d-complete-root-support-and-coefficient-bounds}

Every surviving vacuum contraction in \(W\) is supported on one edge or
the edges of a triangle. The following list accounts for all terms in
(2.7):

\begin{itemize}
\tightlist
\item
  Born-Huang derivatives on two different pair vacua are orthogonal
  after removing the scalar Berry connection. Each remaining term is
  bounded by \(C_{N}/r_{e}^{2}\).
\item
  \(\rho _{0}\) is a sum of pair gradients. Their inner product vanishes
  for disjoint endpoint sets. A cross product on two incident edges has
  weight \(1/(r_{e} r_{f})\) and belongs to their triangle. The
  divergence is pairwise.
\item
  Wick contractions in \(V_{4}\) involve either one repeated root or two
  roots sharing an endpoint. Their weights are \(r_{e}^{-2}\) or
  \((r_{e} r_{f})^{-1}\).
\item
  The first term of \(\phi _{2}\) is a sum of \(Y_{e}^{2}/r_{e}^{2}\)
  contractions. Its second term follows a root-space path
  \(e \to  f \to  e\) and has \(Y_{g}^{2}/(r_{e} r_{f})\), where e,f,g
  form a triangle. A fast derivative and the momentum pairing remove the
  \(Y_{g}\) scale; the resulting weights are again \(r_{e}^{-2}\) or
  \((r_{e} r_{f})^{-1}\).
\item
  \(D_{1}\) contains a pair spin term with weight \(r_{e}^{-1}\), and
  triangle terms with weights \(r_{e}^{-1}\) or
  \(r_{e}^{-1} \sqrt{r_{g}/r_{f}}\). Squaring on the vacuum preserves
  graph orthogonality. The triangle inequality reduces the latter square
  to \(C w_{t}\).
\item
  The slow-orbital part of \(D_{2}\) paired with \(D_{0}\) gives
  \(j_{e}=O(\kappa /r_{e})\) on endpoint coordinates. Its derivative is
  \(O(r_{e}^{-2})\) and varies from its centered value by
  \(O(\kappa /r_{e}^{2})\). Derivatives of spectator vacua are
  orthogonal in this compression; the products with \(\rho _{0}\) have
  pair/triangle support.
\item
  In \(D_{0}^{*}F^{-1}L^{*}F^{-1} S_{f}\), the two bosonic factors lie
  on distinct roots, so its vacuum compression is zero.
\item
  In \(D_{0}^{*}F^{-1}L^{*}F^{-1} N_{Y}^{*}p_{F}\), nonzero Wick
  contractions force the second root path to close the same triangle.
  The two possible pairings have weights \(\kappa /(r_{f} r_{g})\) and
  \(\kappa /(r_{e} r_{f})\).
\item
  In the two-L term of \(D_{2}\) the same closure gives the preceding
  weights with \(\kappa\) squared.
\item
  In the \(J_{2}\) and \(KK^{*}\) terms, the two orbital center maps
  vanish against each other for disjoint endpoint sets. Surviving
  weights are \(\kappa ^{2}/(r_{e} r_{f})\), with the repeated-edge case
  included.
\item
  The \(\rho _{F2}\) part of \(D_{2}\) has the density supports already
  listed and a prefactor bounded by \(\kappa\) squared after pairing
  with \(D_{0}\).
\end{itemize}

For estimates on finite excited states, a root path need not close under
Wick contraction. It still has a uniform inverse-square coefficient
bound. For example a two-L path has normalized weight

% DISPLAY d.25
% SOURCE sqrt(r_h)/(sqrt(r_e) r_f sqrt(r_g r_j)).
\begin{gather*}
\begin{aligned}
& \sqrt{r_{h}}/(\sqrt{r_{e}} r_{f} \sqrt{r_{g} r_{j}}).
\end{aligned}\notag
\end{gather*}

The triangle relation \(r_{h}\le r_{f}+r_{j}\) bounds this by a sum of

% DISPLAY d.26
% SOURCE 1/sqrt(r_e r_f r_g r_j),
% SOURCE 1/(sqrt(r_e) r_f sqrt(r_g)),
\begin{gather*}
\begin{aligned}
& 1/\sqrt{r_{e} r_{f} r_{g} r_{j}}, \\
& 1/(\sqrt{r_{e}} r_{f} \sqrt{r_{g}}),
\end{aligned}\notag
\end{gather*}

both bounded by a numerical multiple of the sum of inverse-square edge
weights by weighted arithmetic-geometric mean. This observation is
useful for mixed source terms and prevents an illicit global-frequency
comparison.

All of these are finite coefficient expressions in pair matrices with
external gaps comparable to their own \(r_{e}\). Normalize each
occurrence of an internal matrix by the incident \(r_{e}\) in that
factor. Its norm is \(O(\kappa )\). Functional calculus and its required
derivatives have uniform convergent resolvent bounds on that fixed pair
ball. Termwise product estimates in the above inventory therefore prove

% DISPLAY d.27
% SOURCE ||W(A)|| <= C_N W_2,
% SOURCE ||W(A)-W(0)|| <= C_N kappa W_2,
% SOURCE W_2=sum_e r_e^-2.                                      (4.2)
\begin{gather*}
\begin{aligned}
& \Vert W(A)\Vert  \le  C_{N} W_{2}, \\
& \Vert W(A)-W(0)\Vert  \le  C_{N} \kappa  W_{2}, \\
& W_{2}=\sum _{e} r_{e}^{-2}.
\end{aligned}\tag{4.2}
\end{gather*}

The derivative of \(j\) is included in the second estimate. It is not
estimated merely by the small value of \(j\), which would miss a nonzero
centered internal derivative.

\subsection{5. Actual shifted-fast correction and leading operator
defect}\label{d-actual-shifted-fast-correction-and-leading-operator-defect}

On the pair source space the adapted excitation energy is at least
\(c r_{e}\); on the triangle source space it is at least \(c s_{t}\).
The finite matrix resolvent identity, applied to the actual sources
(3.1), gives

% DISPLAY d.28
% SOURCE ||S_e*H_exc^-1 S_e-S_e(0)*H_f(0)^-1 S_e(0)||
% SOURCE       <= C_N kappa r_e^-2,
% SOURCE 
% SOURCE ||S_t*H_exc^-1 S_t-S_t(0)*H_f(0)^-1 S_t(0)||
% SOURCE       <= C_N kappa w_t.                                (5.1)
\begin{gather*}
\begin{aligned}
& \Vert S_{e}^{*}H_{\mathrm{exc}}^{-1} S_{e}-S_{e}(0)^{*}H_{f}(0)^{-1} S_{e}(0)\Vert \\
& \le  C_{N} \kappa  r_{e}^{-2},
\end{aligned}\notag \\
\begin{aligned}
& \Vert S_{t}^{*}H_{\mathrm{exc}}^{-1} S_{t}-S_{t}(0)^{*}H_{f}(0)^{-1} S_{t}(0)\Vert \\
& \le  C_{N} \kappa  w_{t}.
\end{aligned}\tag{5.1}
\end{gather*}

Cross terms between distinct graph spaces vanish. The centered source
lemma, including all internal vacuum derivatives and compression jets,
gives

% DISPLAY d.29
% SOURCE W(0)=S(0)*H_f(0)^-1 S(0).
\begin{gather*}
\begin{aligned}
& W(0)=S(0)^{*}H_{f}(0)^{-1} S(0).
\end{aligned}\notag
\end{gather*}

Consequently the actual leading operator defect satisfies

% DISPLAY d.30
% SOURCE ||W(A)-S(A)*H_exc(A)^-1 S(A)||
% SOURCE       <= C_N kappa W_2.                                (5.2)
\begin{gather*}
\begin{aligned}
& \Vert W(A)-S(A)^{*}H_{\mathrm{exc}}(A)^{-1} S(A)\Vert \\
& \le  C_{N} \kappa  W_{2}.
\end{aligned}\tag{5.2}
\end{gather*}

This is an operator-norm estimate on the complete internal Clifford
module. Its actual fast dressing is

% DISPLAY d.31
% SOURCE Z(A)=-H_exc(A)^-1 S(A),                                 (5.3)
\begin{gather*}
\begin{aligned}
& Z(A)=-H_{\mathrm{exc}}(A)^{-1} S(A),
\end{aligned}\tag{5.3}
\end{gather*}

and obeys the root-graph bounds and spectator-derivative correction from
Appendix \(C\). No interacting internal Hamiltonian occurs in (5.3). The
fast energy E0 remains a separate multiplier with its proved
commutator-square payment.

\subsection{6. All-vector use and the remaining reserve
requirement}\label{d-all-vector-use-and-the-remaining-reserve-requirement}

The graph-local lemma now applies to the actual dressing (5.3): its
coefficients are the explicit graph-local functions above. Smooth pair
functional calculus gives active derivative weights \(C r_{e}^{-5}\) or
\(C(\sum _{t} r_{e}^{-1})^{2}w_{t}/s_{t}\); spectator derivatives add
\(C m_{g} W_{2}\). Therefore

% DISPLAY d.32
% SOURCE H_s[Zf] <= C_N r_*^-3 H_s[f]
% SOURCE              +C_N kappa W_2[f]+C_N W_5[f].              (6.1)
\begin{gather*}
\begin{aligned}
& H_{s}[Zf] \le  C_{N} r_{*}^{-3} H_{s}[f] \\
& +C_{N} \kappa  W_{2}[f]+C_{N} W_{5}[f].
\end{aligned}\tag{6.1}
\end{gather*}

In fact triangle sources have no internal Clifford generator, so outside
derivative commutators are the only issue for their internal potential
charges. The more general graph bound is valid without using this
improvement. All internal momenta in (6.1) are unrestricted.

The identity \(H_{\mathrm{exc}} Z+S=0\) cancels the actual leading
complementary source exactly. The complete internal/center form pairs
involving \(Z\) can be controlled by Cauchy-Schwarz in \(H_{s}\), using
(6.1). The remaining derivative source from \(U\) is incident-edge
weighted and has the same payment as Section \(6\) of the geometric
lemma. These steps are valid form estimates on arbitrary complementary
vectors in the corresponding control norm; they are not a finite
oscillator truncation of those vectors.

What is not supplied here is an inequality extracting that control norm
from the original complementary form while retaining the same Gauss
square. In particular, the expression

% DISPLAY d.33
% SOURCE H_s[v]+H_exc[v]+||D_full v||²
\begin{gather*}
\begin{aligned}
& H_{s}[v]+H_{\mathrm{exc}}[v]+\Vert D_{\mathrm{full}} v\Vert ^{2}
\end{aligned}\notag
\end{gather*}

cannot be added to the right side of the original Hamiltonian merely
because its three terms are nonnegative. \(H_{\mathrm{exc}}\) already
uses the frozen piece of \(D_{\mathrm{full}}\). The desired actual
all-vector comparison still needs a non-double-counted reserve,
compatible with unrestricted slow momenta, and a complete
mixed/remainder estimate in that reserve. Finite coefficient
cancellation (5.2) cannot prove this coercivity.

Thus Sections 2--5 identify and estimate the actual leading H1/H2 source
coefficients. Section \(6\) supplies their internal-momentum payment and
states the precise remaining all-vector obstruction rather than silently
assuming it away.
